%=====================================================================
% 两周时间表 2026-11-09 → 11-22
% 编译：xelatex schedule.tex（两遍）
% 注：theme-violet 的 \callout / \warnbox 是「命令」不是环境，用 \callout{...}
%=====================================================================
\documentclass[11pt]{ctexart}
%=====================================================================
% theme-violet.tex —— 讲义精讲 · 紫花主题
% 用法：\documentclass[11pt]{ctexart}
%       \input{assets/theme-violet}
%       \begin{document} ... \end{document}
% 只用 ctexart 自带能力 + xcolor/enumitem/fancyhdr/pifont/amsmath，
% 不依赖 tikz/tcolorbox，保证 MiKTeX 可直接编译。
%=====================================================================

\usepackage[a4paper,left=2.0cm,right=2.0cm,top=2.1cm,bottom=2.3cm,%
            headsep=12pt,footskip=22pt]{geometry}
\usepackage{amsmath,amssymb}
\usepackage{needspace}
\usepackage{xcolor}
\usepackage{array}
\usepackage{enumitem}
\usepackage{pifont}
\usepackage{fancyhdr}
\usepackage[hidelinks]{hyperref}

%---------------------------------------------------------------- 配色
\definecolor{violet}{RGB}{107,78,143}      % 主色：紫
\definecolor{violetdeep}{RGB}{63,44,92}    % 深紫：标题/强调
\definecolor{petal}{RGB}{183,159,208}      % 花瓣：装饰线、花饰
\definecolor{bandviolet}{RGB}{247,243,251} % 浅紫底：档案/提示框
\definecolor{bandcream}{RGB}{252,249,242}  % 米白底：易错/注意框
\definecolor{rulepurple}{RGB}{217,207,230} % 分隔线
\definecolor{inkgray}{RGB}{92,96,110}      % 次要文字

%---------------------------------------------------------------- 字体
% 有 Noto CJK 用 Noto，没有就退回 ctex 默认（Windows 下为宋体/黑体）
\IfFontExistsTF{Noto Serif CJK SC}{\setCJKmainfont{Noto Serif CJK SC}}{}
\IfFontExistsTF{Noto Sans CJK SC}{\setCJKsansfont{Noto Sans CJK SC}}{}

%---------------------------------------------------------------- 版式
\setlist{nosep,leftmargin=1.9em}
\linespread{1.28}
\setlength{\parskip}{3.5pt}
\setlength{\parindent}{0pt}
\setlength{\tabcolsep}{5pt}
\renewcommand{\arraystretch}{1.28}

%---------------------------------------------------------------- 花饰
% \ding{95} = 五瓣小花，\ding{96} = 星形花；都来自 pifont，必定存在
\newcommand{\orn}{\textcolor{petal}{\ding{95}}}          % 行内小标
\newcommand{\ornstar}{\textcolor{petal}{\ding{96}}}      % 星形花
\newcommand{\flowerrule}{%                                % 居中花饰分隔线
  \par\vspace{7pt}%
  \noindent\makebox[\linewidth]{\color{petal}%
    \rule{0.20\linewidth}{0.6pt}\hspace{8pt}{\small\ding{95}}\hspace{4pt}%
    \raisebox{1.1pt}{\tiny\ding{95}}\hspace{4pt}{\small\ding{95}}\hspace{8pt}%
    \rule{0.20\linewidth}{0.6pt}}%
  \par\vspace{7pt}}

%---------------------------------------------------------------- 标题与节头
\newcommand{\doctitle}[3]{%   #1 主标题  #2 副标题  #3 元信息（可用 \\ 分行）
  \begin{center}
    {\color{petal}\small \orn\enspace #2\enspace \ornstar}\\[7pt]
    {\LARGE\bfseries\color{violetdeep} #1}\\[10pt]
    {\color{petal}\rule{0.55\linewidth}{1pt}}\\[4pt]
    {\small\color{inkgray} #3}
  \end{center}\vspace{2pt}}

\newcommand{\hdr}[2]{%   #1 章节号（如 §0）  #2 标题
  \par\vspace{12pt}%
  \noindent\colorbox{violetdeep}{\color{white}\bfseries\small\;#1\;}\hspace{7pt}%
  {\color{violetdeep}\large\bfseries #2}\par\nobreak
  \vspace{3pt}{\color{petal}\hrule height 1.2pt}\vspace{6pt}}

\newcommand{\subhdr}[1]{%
  \par\vspace{6pt}%
  \noindent{\color{violet}\bfseries #1}\par\vspace{1pt}}

%---------------------------------------------------------------- 提示框
\newcommand{\callout}[1]{% 浅紫底：讲义档案、读法说明
  \par\medskip\noindent
  \fcolorbox{rulepurple}{bandviolet}{\parbox{\dimexpr\linewidth-2\fboxsep-2\fboxrule}{\vspace{2pt}#1\vspace{2pt}}}%
  \par\medskip}

\newcommand{\warnbox}[1]{% 米白底：易错、边界条件
  \par\medskip\noindent
  \fcolorbox{rulepurple}{bandcream}{\parbox{\dimexpr\linewidth-2\fboxsep-2\fboxrule}{\vspace{2pt}#1\vspace{2pt}}}%
  \par\medskip}

\newcommand{\gapitem}[2]{% #1 缺口类型（紫） #2 内容
  \par\vspace{2pt}\noindent
  \hangindent=1.4em\hangafter=1
  {\color{petal}\small\ding{96}}\enspace{\bfseries\color{violetdeep}#1}\enspace #2\par}

% 题干块（\qstem）：每道例题／作业的「详解」开头必须先复述题干与出处。
% 起因：2026-10-25 第五章讲义把七道作业的题面集中放在一页、详解另起一页，
% 读者翻到「第 4 题」时看不到条件，等于一道没有题干的题目。
% 规则：详解块内自带题干，读者不翻页就能读懂题在问什么。
\newcommand{\qstem}[1]{%
  \par\medskip\noindent
  \fcolorbox{violet}{bandviolet}{%
    \parbox{\dimexpr\linewidth-2\fboxsep-2\fboxrule}{%
      \vspace{2pt}%
      {\small\bfseries\color{violet}题干}\hspace{0.8em}#1%
      \vspace{2pt}}}%
  \par\medskip}

% v2: 保护题干与开头，长题可跨页，续页补上下文。
\newcommand{\qbreak}{\needspace{6\baselineskip}}
% 独立练习：ID、来源、题干、考点、答案标签（通常 ans:ID）。
\newcommand{\practice}[5]{%
  \qbreak\subhdr{练习 #1}\qstem{（#2）#3}%
  \textbf{考点}：#4\par
  详解见附录第~\pageref{#5}~页。\par\medskip}
% 附录开头：ID、来源、条件摘要；自动产生 ans:ID 标签。
\newcommand{\answerstart}[3]{%
  \qbreak\subhdr{#1\quad 详解}\label{ans:#1}\qstem{（#2）#3}}
\newcommand{\answercontinue}[2]{%
  \needspace{5\baselineskip}\subhdr{#1\quad 续解}%
  {\small\color{inkgray}本页接续：#2}\par\medskip}
% 延迟提供默认值，让旧讲义可在导言区保留自己的 \newcommand{\step}。
\AtBeginDocument{\providecommand{\step}[1]{\needspace{3\baselineskip}\par\medskip\textbf{#1}\par}}

% 前提旁注：只放在「需要理解的那一步」后面，给陈述 + 本题代入，不重证。
\newcommand{\aside}[1]{%
  \par\vspace{1pt}\noindent
  \hangindent=1.6em\hangafter=1
  {\small\color{inkgray}%
    \hspace*{0.6em}{\color{petal}$\triangleright$}\hspace{2pt}#1}%
  \par\vspace{2pt}}

%---------------------------------------------------------------- 表格
% 本机 MiKTeX 的 colortbl 与 array 版本不匹配（\rowcolor 会报错），
% 故表格统一用「上粗紫线 + 表头 + 中细线」的三线式，不依赖 colortbl。
\newcommand{\tabtop}{\par\vspace{3pt}{\color{violet}\hrule height 1.2pt}\vspace{1pt}}
\newcommand{\tabmid}{\par\vspace{1pt}{\color{petal}\hrule height 0.5pt}\vspace{1pt}}
\newcommand{\tabbot}{\par\vspace{3pt}{\color{rulepurple}\hrule height 0.5pt}\vspace{8pt}}
\newcommand{\thead}[1]{{\bfseries\color{violetdeep}#1}}

%---------------------------------------------------------------- 页眉页脚
\pagestyle{fancy}
\fancyhf{}
\renewcommand{\headrulewidth}{0.4pt}
\renewcommand{\headrule}{\hbox to\headwidth{\color{rulepurple}\leaders\hrule height \headrulewidth\hfill}}
\fancyhead[L]{\small\color{inkgray}\songti \doctag}
\fancyhead[R]{\small\color{inkgray}\songti \coursetag}
\fancyfoot[C]{\small\color{inkgray}\orn\enspace{\color{violet}\thepage}\enspace\orn}

% 每篇讲义在导言区设这两个值
\newcommand{\doctag}{}
\newcommand{\coursetag}{}
\newcommand{\settags}[2]{\renewcommand{\doctag}{#1}\renewcommand{\coursetag}{#2}}

%---------------------------------------------------------------- 数学记号
\newcommand{\vect}[1]{\boldsymbol{#1}}
\newcommand{\dd}{\,\mathrm{d}}
\newcommand{\ddx}{\frac{\mathrm{d}}{\mathrm{d}x}}
\addtolength{\jot}{2pt}




\settags{两周时间表 · 11-09 → 11-22}{111 学习助手 · 中国时间}
\setmainfont{Times New Roman}[Ligatures=TeX]

\renewcommand{\callout}[1]{%
  \par\medskip\noindent
  
  \fcolorbox{rulepurple}{bandviolet}{\parbox{\dimexpr\linewidth-2\fboxsep-2\fboxrule\relax}{\vspace*{5pt}#1\vspace*{5pt}}}%
  \par\medskip}
\renewcommand{\warnbox}[1]{%
  \par\medskip\noindent
  
  \fcolorbox{rulepurple}{bandcream}{\parbox{\dimexpr\linewidth-2\fboxsep-2\fboxrule\relax}{\vspace*{5pt}#1\vspace*{5pt}}}%
  \par\medskip}

\newcommand{\sech}[2]{%
  \par\vspace{13pt}\needspace{5.5cm}%
  {\color{violetdeep}\bfseries\large #1#2}\par
  \vspace{2pt}{\color{petal}\hrule height 1.2pt}
  \vspace{7pt}}

\newcommand{\dayblock}[4]{%
  \par\vspace{6pt}\needspace{3.0cm}%
  \noindent{\colorbox{violetdeep}{\color{white}\bfseries\small\hspace{0.45em}#1\hspace{0.45em}}}\hspace{6pt}%
  {\color{violetdeep}\bfseries #2}\hfill{\small\color{violet}主任务：#3}\par
  \vspace{2pt}{\color{rulepurple}\hrule height 0.5pt}\vspace{3pt}%
  \renewcommand{\arraystretch}{1.12}\begin{tabular}{@{}p{1.05cm}p{\dimexpr\linewidth-1.45cm\relax}@{}}
  #4
  \end{tabular}\par}

\newcommand{\slota}[1]{\textbf{\color{violet}上午} & #1\\[2pt]}
\newcommand{\slotp}[1]{\textbf{\color{violet}下午} & #1\\[2pt]}
\newcommand{\slote}[1]{\textbf{\color{violet}晚上} & #1\\[1pt]}

\begin{document}

\doctitle{未来两周时间表}{2026-11-09（周一）→ 11-22（周日）· 第 11–12 周}%
{中国时间 UTC+8\enspace·\enspace 8 个固定节点 + 1 场竞赛 + 4 次实验 + 2 场考试 + 1 个口试\\
生成于 2026-11-09（周一）\enspace·\enspace 时间均为到场或开始时间，不含路上}

\flowerrule

\sech{一、}{这两周只认三条主线}

\callout{%
\textbf{本周（11-09 → 11-13）：俄语。}周五 15:35 考试，之前每一天都有一块俄语，别的一天不排新东西。\\
\textbf{下周（11-15 → 11-20）：微分方程。}周六 10:05 考试，17 条考纲全考（见第四节），按讲义一份一份过。\\
\textbf{贯穿两周：马原免试口试。}第 12 周（11-16 至 11-22）内某天，日期还没定；所以把两个「收口日」压在 11-15 一天里做完。\\[3pt]
\textbf{其它课一律只跟课、不推进}：理力、光学、数物二、电工（实验除外）这两周不做额外复习。\\
\textbf{每天两件固定动作}：睡前 10 分钟马原口述 3 个镜头（11-10 起）；22:30 收工，不熬。}

\sech{二、}{固定节点（迟到不起的那种）}

\tabtop
{\footnotesize\setlength{\tabcolsep}{3pt}
\begin{tabular}{@{}p{1.9cm}p{2.3cm}p{5.75cm}p{3.05cm}p{2.7cm}@{}}
\thead{日期} & \thead{时间} & \thead{事项} & \thead{地点} & \thead{备注}\\
\hline
11-10 二 & 10:00–12:15 & 电工实验 6：集成运放（\textbf{25 分}考试实验） & 综一 317 & 场次待确认 \\
11-13 五 & 15:35–17:15 & 大学俄语 1 · 结课考试 & 综 353 & 已确认 \\
11-14 六 & 09:00–11:30 & 全国大学生数学竞赛（预赛） & 待确认 & 准考证自打 \\
11-14 六 & 13:30–17:30 & 大物实验 3.12／3.13 分光计、光栅常量 & 综合楼 308 & 已预约 \\
11-16 至 22 & 未定 & 马原免试口试（第 12 周） & 未定 & 日期待确认 \\
11-19 四 & 13:00–15:15 & 电工实验 7：集成逻辑门 & 综一 317 & 已预约 \\
11-21 六 & 10:05–11:45 & 微分方程 · 期末考试 & 综 160 & 17 条考纲 \\
11-21 六 & 13:30–15:30 & 大物实验 5.3 温度计设计 & 综合楼 205 & 考完直接去 \\
\end{tabular}}
\tabbot

\callout{\textbf{每天真正空着的时段}（按课表扣掉上课）：周一 15:35 后；周二 10:00 后；周三 10:00–11:40 与 21:20 后；周四 13:30 后；周五 13:30 后；周末全天（11-14 除外）。下面表里排的任务就落在这几个格子里。}

\sech{三、}{逐日安排}

\dayblock{11-09}{周一 · 第 11 周（今天）}{电工实验 6 预习}{%
\slota{08:00 理力（综113）\enspace｜\enspace 10:00 马原（综309）}
\slotp{13:30 数物二（综251）· 第七章开头}
\slote{20:00–21:40 电工实验 6 预习：实验6·手册 p19–22 逐页对照执行单 ＋ 备考讲义里的评分点\newline 21:40–22:10 俄语 30 分：形容词词尾卡}}

\dayblock{11-10}{周二}{俄语变格与词性}{%
\slota{08:00 光学（一馆1-212）\enspace｜\enspace 10:00–12:15 电工实验 6（综一317）}
\slotp{14:00–15:30 俄语：变格表 ＋ УРОК1–8 总复习课件前半（一卡一题，遮答案自测）}
\slote{20:00–22:00 俄语：全课程精讲 §2 名词 ＋ §4 形容词（只读这两节，读完各做卡上 3 题）}}

\dayblock{11-11}{周三}{俄语词性与复数}{%
\slota{08:00 马原（综309）\enspace｜\enspace 10:00 起 俄语 60 分：-ь 词性 38 词卡 ＋ 复数三表}
\slotp{13:30 形势与政策（综B501）\enspace｜\enspace 15:35 电工（综213）}
\slote{19:40–21:20 改革开放史（综153）\enspace｜\enspace 21:20 起 俄语 40 分：总复习课件后半}}

\dayblock{11-12}{周四}{俄语错题＋打印}{%
\slota{08:00 理力（综113）\enspace｜\enspace 10:00 体育（体育馆）}
\slotp{13:30–15:30 俄语：只过错题（形容词词尾、词性判断）；\textbf{15:30–16:00 打印准考证与实验数据表}}
\slote{20:00–21:30 俄语：УРОК8 双语教程快速过 ＋ 阅读限时一篇}}

\dayblock{11-13}{周五}{俄语考试}{%
\slota{08:00 电工（综213）\enspace｜\enspace 10:00 光学（一馆1-212）}
\slotp{13:30–15:00 只看知识卡，不碰新题（精讲 §8 四张速查表）\enspace｜\enspace \textbf{15:35–17:15 俄语考试（综353）}}
\slote{不做马原。收拾明天：准考证、身份证、黑色签字笔、已打印的实验数据表、午饭地点}}

\dayblock{11-14}{周六}{竞赛 ＋ 实验（当天不排学习）}{%
\slota{08:00 出发（考场见准考证）\enspace｜\enspace 09:00–11:30 数学竞赛预赛}
\slotp{13:30–17:30 大物实验 3.12／3.13：分光计调节 ＋ 光栅常量（综合楼308）}
\slote{休息。21:00 前把马原逐字稿、PPT、手卡材料放到桌上，明天上午直接用}}

\dayblock{11-15}{周日}{马原收口（上午一场、晚上一场）}{%
\slota{09:00–12:00 收口日 1：全文朗读 5 遍、改 2 处、PPT 定稿、打印手卡}
\slotp{14:00–16:00 微分方程：L13 方程组入门 ＋ L14–15 常系数齐次组（两份精讲，读后各做 2 题）}
\slote{20:00–22:30 收口日 2：计时演练 4 遍（三遍落进 7 分 20 秒 – 7 分 40 秒）＋ 手机录像 1 遍}}

\dayblock{11-16}{周一 · 第 12 周}{微分方程矩阵指数}{%
\slota{08:00 理力（综113，本周最后一次）\enspace｜\enspace 10:00 马原（综309）}
\slotp{13:30 数物二（综251）}
\slote{20:00–20:40 马原全套 1 遍 ＋ 抽纸条 2 题\enspace｜\enspace 20:40–22:30 微分方程 L16–17 矩阵指数}}

\dayblock{11-17}{周二}{微分方程非齐次组}{%
\slota{08:00 光学（一馆1-212）\enspace｜\enspace 10:00–12:00 微分方程 L17–18 非齐次组（参数变易 ＋ 待定系数）}
\slotp{14:00–15:30 微分方程：L16–18 的例题各闭卷做 1 题}
\slote{20:00–20:40 马原 1 遍\enspace｜\enspace 20:40–22:30 微分方程 错题回炉}}

\dayblock{11-18}{周三}{微分方程相图与稳定性}{%
\slota{08:00 马原（综309）\enspace｜\enspace 10:00–11:40 微分方程 L21–23 临界点与稳定性}
\slotp{13:30 形势与政策（综B501）\enspace｜\enspace 15:35 电工（综213）}
\slote{19:40 改革开放史（综153）\enspace｜\enspace 21:30–22:15 电工实验 7 预习}}

\dayblock{11-19}{周四}{实验 7 ＋ 微分方程限时自测}{%
\slota{08:00 理力（综113，结课）\enspace｜\enspace 10:00 体育（体育馆）}
\slotp{13:00–15:15 电工实验 7（综一317）\enspace｜\enspace 15:30–17:00 微分方程：期中之后总纲过一遍}
\slote{20:00–20:40 马原 1 遍\enspace｜\enspace 20:40–22:30 微分方程 L1–L12 乱序自测（限时 90 分）}}

\dayblock{11-20}{周五}{微分方程考前收口}{%
\slota{08:00 电工（综213）\enspace｜\enspace 10:00 光学（一馆1-212）}
\slotp{13:30–15:00 微分方程：只过错题 ＋ L1–L12 考点自查清单}
\slote{20:00–21:00 微分方程最后过一遍知识卡\enspace｜\enspace 22:30 睡}}

\dayblock{11-21}{周六}{考试 ＋ 实验}{%
\slota{09:20 到综160\enspace｜\enspace 10:05–11:45 微分方程考试}
\slotp{13:30–15:30 大物实验 5.3 温度计设计（综合楼205）}
\slote{休息，不安排任何学习}}

\dayblock{11-22}{周日}{报告清账 ＋ 理力恢复}{%
\slota{09:00–11:30 补实验报告：3.12／3.13 分光计 ＋ 5.3 温度计}
\slotp{14:00–16:30 理力 T5 封口：能量积分、转向点、有效势（T05 精讲）＋ T6 碰撞（T06）}
\slote{20:00–21:00 马原口试彩排 1 遍\enspace｜\enspace 21:00–22:00 排下周（11-23 起）}}

\sech{四、}{微分方程 11-21 考什么（课件作者 11-05 发的 17 条）}

\warnbox{%
\textbf{一阶}\enspace 定义与初值问题；可分离、齐次、可化齐次；恰当方程与积分因子；一阶线性、Riccati、Bernoulli；齐次极坐标方程；关于 $y$ 或关于 $x$ 可解、Clairaut、Lagrange。\\
\textbf{高阶}\enspace 降阶（含不含 $x$、不含 $y$ 或前几阶导数）；线性相关与 Wronskian、基础解系；高阶线性（齐次／非齐次、参数变易、常系数）；Cauchy–Euler；幂级数法、Frobenius。\\
\textbf{方程组}\enspace 定义与 Wronskian、$n$ 阶方程化组；常系数齐次组；$e^{Ax}$；非齐次组（参数变易／待定系数）。\\
\textbf{定性}\enspace 齐次线性系统相图、临界点类型与稳定性；一阶近似下的稳定性。\\[2pt]
对应讲义：L13、L14–15、L16–17、L17–18、L18–L20、L21–L23 精讲 ＋《期中之后·全景总纲》＋ L1–L12 乱序自测。}

\sech{五、}{还要你确认的 6 件事}

\begin{enumerate}[label=\textbf{\arabic*.},leftmargin=1.6em]
\item \textbf{竞赛考场在哪、几点要求进场}（准考证／考生须知）——决定 11-14 从 11:30 到 13:30 够不够赶到综合楼。
\item \textbf{电工实验 6 你是哪一场}：10:00–12:15（2502 班 20 人）还是 15:30–17:45（两班各 9 人）。
\item \textbf{大物 3.12／3.13 与 5.3 的讲义、预习页}——库里只有 3.8 和 3.11 两份，没有底本就做不出讲义和数据表。
\item \textbf{马原口试的确切日期与材料截止}——它决定 11-15 这天能不能安心让给收口。
\item \textbf{数物二是不是只有周一 5–6 节}（第 6 周截图里周二也出现过一次）。
\item \textbf{理力、光学的考试日期}（教务系统还没出）。
\end{enumerate}

\sech{六、}{这两周之外已经埋着的三处}

\callout{%
\textbf{11-24（周二）} 电工实验 8：触发器与计数器，10:00–12:15。\\
\textbf{11-29（周日）上午硬冲突}：影视鉴赏考试 09:10–10:10 与 大物实验 4.9 光纤 08:00–10:00 重叠 —— 现在就要找老师换场，别拖到那周。\\
\textbf{12-06（周日）} 马原机考 08:00–08:40 结束，09:00 电工综合设计实验二开始，只有 20 分钟转场。}

\flowerrule
{\footnotesize\color{inkgray}
依据：课表（节次时间与周次）、考试日程（考试与考场）、大物实验六场安排、电工实验八次排课、
马原免试口试作战方案（收口日与里程碑）、微分方程课件作者 11-05 发的 17 条考纲。\\
本表是排期建议，不是已确认事实：标「待确认」的六项以你的回复为准，改一版重发即可。}

\end{document}

